\documentclass{ximera}

\input{../../preamble.tex}

\title{Euler's Formula}

\begin{document}

\begin{abstract}
complex exponentials
\end{abstract}
\maketitle






\subsection*{An Adventure}




One of the biggest surprises in mathematics is the fact that trigonometric functions and exponential functions are the same building blocks, provided we use Complex numbers. 


Euler's Formula is the bridge:



\[    e^{i \, t} = \cos(t) + i \sin(t)    \]



This is very weird.  It claims that exponential functions are trigonometric functions are two perspectives on the same thing.  Our previous mathematcis courses have never even hinted at such a connection.  That is because the bridge  travels through the Complex Numbers.



With this bridge, we now have direct connections between


\begin{itemize}
\item percentage growth
\item exponential functions
\item trigonometric functions
\item triangles
\item circles
\end{itemize}

This would also connect inverse trigonometric functions and logarithms and soon hyperbolas and hyperbolic trigonometric functions.  This is crazy!


\begin{center}
\textbf{\textcolor{red!80!black}{Everything is Connected!!!}}
\end{center}






It would be nice to show that Euler's Formula is true.  However, we will need Calculus to dot our i's and cross our t's and leave no doubt.   


But, we can understand the basic idea. 


We just have to remember properties of $\cos(t) + i \sin(t)$ that we have seen.




\textbf{$\blacktriangleright$ Multiplication}


$\cos(t) + i \sin(t)$ is a Complex number on the unit circle and we have seen that multiplication is accomplished by adding the angles.



\[ (\cos(\theta_1) + i \sin(\theta_1)) \cdot (\cos(\theta_2) + i \sin(\theta_2))  = \cos(\theta_1 + \theta_2) + i \sin(\theta_1 + \theta_2) \]


Just like $A^B \cdot A^C = A^{B+C}$




\textbf{$\blacktriangleright$ Quotients}


$\cos(t) + i \sin(t)$ is a Complex number on the unit circle and we have seen that division is accomplished by subtracting the angles.


\[  \frac{\cos(\theta_1) + i \sin(\theta_1)}{\cos(\theta_2) + i \sin(\theta_2)} = \cos(\theta_1 - \theta_2) + i \sin(\theta_1 - \theta_2) \]


Just like $\frac{A^B}{A^C} = A^{B-C}$


The arithmetic of $\cos(t) + i \sin(t)$  is the same as the exponential arithmetic.



To get the exact exponential form, we will need derivative.

You can read the whole story here.



















\begin{center}
\textbf{\textcolor{green!50!black}{ooooo-=-=-=-ooOoo-=-=-=-ooooo}} \\

more examples can be found by following this link\\ \link[More Examples of Complex Bridge]{https://ximera.osu.edu/csccmathematics/precalculus2/precalculus2/complexBridge/examples/exampleList}

\end{center}







\end{document}
